\section{Introduction}
\label{sec:intro}

A two-sample test answers whether two populations are identical. In many applications they are known not to be: the
two arms of an observational study, a training sample and the population on which a model will be deployed, a
reference cohort and a new batch. The useful question is then how different they are, and the contamination model
\citep{huber1964robust} gives it a precise answer. The total-variation distance
\[
\TV(P,Q)=\sup_A\,|P(A)-Q(A)|
\]
is the largest difference between the probabilities that two laws $P$ and $Q$ give to one event. The two laws can be
written as mixtures $P=(1-\alpha)R+\alpha P_0$ and $Q=(1-\alpha)R+\alpha Q_0$ that share a component $R$ of weight
$1-\alpha$ exactly when $\TV(P,Q)\le\alpha$. The smallest such $\alpha$, $\alpha^\star=\TV(P,Q)$, which we call the
contamination level, is therefore the fraction of each population with no counterpart in the other. In an
observational study it is the smallest fraction whose removal from each arm lets the covariate distributions
coincide. We certify lower bounds on $\alpha^\star$ from two samples, of sizes $n$ and $m$, on any sample
space that carries a bounded characteristic kernel, with finite-sample guarantees and without a model for how the
populations differ, and we also return the trimming of each sample at that level.

Discarding part of a population is a trimming. An $\alpha$-trimming of $P$ is a law obtained from $P$ by removing at
most a fraction $\alpha$ of its mass, possibly only partially at each point, and renormalizing, that is, a law with
density at most $(1-\alpha)^{-1}$ relative to $P$ \citep{gordaliza1991best,alvarez2008trimmed}.
\citet{alvarez2012similarity} observed that two laws have a common $\alpha$-trimming exactly when
$\TV(P,Q)\le\alpha$, so $\alpha^\star$ is the smallest level at which $P$ and $Q$ can be trimmed into the same law.
Their test of $\TV(P,Q)\le\alpha$ trims both samples to minimize a Wasserstein distance, with a bootstrap guarantee
that is asymptotic, pointwise and univariate, and they ask whether ``a different choice of metric could lead to some
type of uniform bound'' \citep[p.~617]{alvarez2012similarity}. With a
kernel discrepancy the answer is positive at every sample size, which gives a lower confidence bound on $\alpha^\star$.
An upper bound needs assumptions: data can show that $\TV(P,Q)$ is large, not that it is small
(Theorem~\ref{thm:impossible}).

A bounded kernel $k$ on the sample space maps each point $x$ to a feature vector $\phi(x)$ in a Hilbert space $\Hh$,
the reproducing kernel Hilbert space (RKHS) of $k$, and each law $P$ to the mean of its feature vectors, the kernel
mean embedding $\mu_P=\E_P[\phi(X)]$. The distance $\|\mu_P-\mu_Q\|_\Hh$ is the maximum mean discrepancy (MMD) of
\citet{gretton2012kernel}, and for a characteristic kernel, such as the Gaussian or the Laplace kernel, the embedding
determines the law. Our approach is geometric: instead of an optimal trimming, we follow the set of everything the
trimmings can produce. The trimmed region $\mathcal C_\alpha(P)$, the set of the embeddings of all $\alpha$-trimmings
of $P$, is a convex compact subset of $\Hh$ that grows with $\alpha$ from the single point $\mu_P$, and the regions of
$P$ and $Q$ first meet exactly at $\alpha^\star$, since a point common to both is the embedding of a trimming of $P$
and of a trimming of $Q$, which must then be the same law (Figure~\ref{fig:regions}). The distance $D_\alpha(P,Q)$
between the regions is thus positive below $\alpha^\star$ and zero from $\alpha^\star$ on. For the empirical laws the
regions are polytopes, whose distance is one convex quadratic program, and to show that $\TV(P,Q)>\alpha$ is to show
that two compact convex sets are apart, which the direction that joins their closest points shows.

The difference of the two closest points is the witness, a function on the sample space that is large where the first
trimmed law has more mass than the second. Once the witness is fixed, choosing the trimming is a linear problem of
bathtub type, so the optimal trimmings keep at full weight the observations on one side of a cut of the witness and
discard those beyond it, only the observations at the cut sharing weight. The quadratic program that gives the
distance, solved on all the data at the certified level, therefore also gives the trimmed samples. They converge to
the best trimming at a fixed level, at the rate $(n\wedge m)^{-1/2}$ when the witness has a density at the cuts, and to
the common component of $P$ and $Q$ at the estimated level, the target of a cohort of comparable patients.

Any function $\rp$ with values in $[0,1]$ has $\E_P\rp-\E_Q\rp\le\TV(P,Q)$, so a test that shows
$\E_P\rp-\E_Q\rp>\alpha$ shows that $\TV(P,Q)>\alpha$, and the geometry says which $\rp$ to use. In a direction
$f\in\Hh$ a region reaches as far as a trimmed mean, the mean of $f(X)$ over the upper $(1-\alpha)$-fraction of its
values, so the separation of the two regions in that direction is a difference of two trimmed means. A trimmed mean is
an infimum of ordinary means, $\inf_t\{t+(1-\alpha)^{-1}\E_P[(f(X)-t)_+]\}$, one for each cut $t$
\citep{rockafellaruryasev2000}. At the two cuts at which the trimmings at the tested level cut, the separation in the
direction of the witness is at most a positive multiple of the excess $\E_P\rp-\E_Q\rp-\alpha$ of a ramp $\rp$ over the
level, where $\rp$ is $0$ below the lower cut, $1$ above the upper one and linear in between, so a positive separation
gives a positive excess. We fit the witness and its cuts on one part of the sample. As the ramp is bounded and fixed
before the rest of the data are seen, betting on it gives a test valid at every sample size with no assumption on the
laws, which rejects the faster the larger the excess. Our
lower confidence bound $\widehat\alpha$ runs this test on a grid of levels fixed in advance, read upward to the first
level not rejected, and its error is measured against the mesh of the grid, the largest gap between its levels.

Betting, like classification, is a problem of estimating a set. For every $[0,1]$-valued $\rp$,
$\E_P\rp-\E_Q\rp=\alpha^\star-\mathrm{err}(\rp)$, with $\mathrm{err}(\rp)$ the distance under $|P-Q|$ from $\rp$ to the
indicator of the set where $P$ has the larger density (Section~\ref{sec:linearmargin}), so betting on $\rp$ certifies
levels near $\alpha^\star$ only if $\rp$ is close to that indicator. The trimming enters this estimate twice. The cut
points of the tested level carry most of the finite-sample gain over a ramp of the untrimmed witness in our simulations
(Section~\ref{sec:numharder}), while refitting the witness at that level makes the bound consistent where a ramp of the
untrimmed witness need not be (Proposition~\ref{prop:loconsistency}, Section~\ref{sec:alphazero}), gives the rates
below and returns the trimming. Any such $\rp$ fitted on earlier data, a classifier for instance, can replace the
ramp, and tuned classifiers certify more where a half-space or a tree is close to that set. A threshold valid under
every pair in the null also turns the empirical distance $\widehat D_\alpha$ into an analytic estimator, which needs no
split but many more observations, our baseline.

The geometry also governs the rates. No valid lower bound has a rate uniform over all pairs
(Section~\ref{sec:upper}), and the rate of ours depends on how fast $D_\alpha$ grows as $\alpha$ decreases below
$\alpha^\star$. A margin $D_{\alpha^\star-\epsilon}(P,Q)\ge c_0\epsilon^\gamma$ guarantees an error of order
$(n\wedge m)^{-1/(2\gamma)}$ up to the mesh. The margin is linear, $\gamma=1$, when a function of the RKHS separates
the two contaminations, the parts of $P$ and $Q$ beyond their common component, and the bound computed on a grid of
mesh $(n\wedge m)^{-1/2}$ is then within order $(n\wedge m)^{-1/2}$ of $\alpha^\star$, which is minimax over such
pairs. When the contaminations touch, as for continuous densities that coincide on no open set, the path
$\alpha\mapsto D_\alpha$ leaves zero tangentially for every continuous kernel, and the margin guarantees only
$(n\wedge m)^{-2/(2s+3)}$ for a Sobolev kernel of order $s$ on the line and a logarithmic rate for the Gaussian kernel,
with no lower bound proved to match.

The margin is pessimistic: it bounds the excess of the population ramp, the ramp of the population witness between
its population cuts, through the gap between these cuts, which it replaces by its largest value. When
$D_\alpha$ is of exact order $(\alpha^\star-\alpha)^\gamma$, the population ramp keeps an excess of at least a fixed
fraction of $\alpha^\star-\alpha$ whatever $\gamma$ (Proposition~\ref{prop:excess}), where the margin sees an excess of
order $(\alpha^\star-\alpha)^\gamma$ only. For the Laplace kernel on the line we show
that the fitted ramp keeps it too: for densities that cross once, under a condition \eqref{eq:condR} on the likelihood
ratio that excludes contaminations local to a common core, the error on a grid of mesh $(n\wedge m)^{-1/2}$ is of order
$(n\wedge m)^{-1/2}$ up to logarithmic factors. A one-sided Kolmogorov--Smirnov bound attains this rate there, and a
bound that does not sort the data, and also returns the set, matches it. The Laplace witness hurts the trimming: it stays
between its cut points, so the observations at the cut tie, while our rate for the trimming at a fixed level assumes a
witness with a density at the cut, as with the Gaussian kernel. So the two kernels divide the work, the Laplace kernel
for the bound and the Gaussian kernel where the trimming needs a rate. In $\R^d$ we keep the Laplace kernel as the
default, because its witness still sees differences of high frequency that the Gaussian witness at the same bandwidth
misses, as on our one-dimensional oscillating densities, although it trails the Gaussian kernel slightly in our
multivariate designs; Section~\ref{sec:discussion} conjectures the rate attainable in $\R^d$.

On the Right Heart Catheterization study of \citet{connors1996rhc}, comparing $\rrRhcNcovFull$ baseline covariates of
the patients who received a catheter and of those who did not, our bound shows that $\TV(P,Q)>\mnRdRhcMin$.
Before the covariate distributions can coincide, one arm, or each if both lose the same fraction, must lose more than
that fraction. This is not what the rules that discard patients by their estimated propensity score control
\citep{crump2009overlap,lei2021overlap,damour2021overlap}: they aim at positivity for a weighted comparison, and a total
variation of this size is compatible with propensity scores bounded well away from $0$ and $1$. On accelerometer curves from the National Health and Nutrition Examination Survey the bound measures how
far apart age and body-mass groups are.

Our bound is an instance of the lower bounds on $\TV$ of \citet{naf2022hplb}, which evaluate a function fitted on
other data, calibrated by the betting test of \citet{shekhar2024nonparametric} with a trimmed null. Such bounds and the
classifier tests \citep{lopezpaz2017revisiting,kim2021classification} apply on any sample space. What the geometry adds
is a canonical function to bet on, fitted at the level tested, for which the bound is consistent and has rates, and
the trimmed samples that come with it. As the Bayes risk of telling $P$ from $Q$ with equal priors is
$\{1-\TV(P,Q)\}/2$, the bound is also an upper confidence bound on a Bayes error, usually estimated
\citep{berisha2016empirically,noshad2019learning,theisen2021evaluating}. \citet{horebarber2026blurred} bound a blurred
relaxation of $\TV$. Mixture proportion estimation \citep{blanchard2010semi,scott2015rate,ramaswamy2016mixture} targets a
one-sided relative, $1-\TV(P,Q)$ being the largest weight of a component common to $P$ and $Q$, with the same dichotomy
of rates (Section~\ref{sec:optimality}), and finite-sample lower bounds on a mixing proportion go back to the
proportion of true null hypotheses in multiple testing
\citep{genovese2004stochastic,meinshausen2006estimating,patra2016estimation}. None of them trims both samples, and
\citet{nafshi2024partial}, when a known fraction of one sample has been assigned to the other, trim the receiving sample
alone. In one sample, the mixture index of fit of \citet{rudas1994index} \citep[see also][]{liu2009tolerance} is the
smallest fraction of the data that must be trimmed for the rest to follow a model, and \citet{delbarrio2020ks} bound it
on the line with the Kolmogorov--Smirnov distance. Other related work is cited where it is used.

Sections~\ref{sec:regions} to~\ref{sec:trimming} treat the regions, the test and the analytic baseline, the
split-sample bound, its rates and the trimming, Section~\ref{sec:numerical} the simulations and the data, and the
supplement gives the proofs and the experimental details. All computations use \texttt{trimmedmmd}, a Python package
released at \url{https://github.com/HristoInouzhe/trimmed-mmd}.
